\documentclass[12pt,a4paper]{article}
\usepackage{cmap} % Makes the PDF copiable. See http://tex.stackexchange.com/a/64198/25761
\usepackage[T1]{fontenc}
\usepackage[brazil]{babel}
\usepackage[utf8]{inputenc}
\usepackage{amsmath}
\usepackage{amsfonts}
\usepackage{amssymb}
\usepackage{amsthm}
\usepackage{textcomp} % \degree
\usepackage{gensymb} % \degree
\usepackage[usenames,svgnames,dvipsnames]{xcolor}
\usepackage{hyperref}
\usepackage{graphicx}
\usepackage[margin=2cm]{geometry}

\hypersetup{
    colorlinks = true,
    allcolors = {blue}
}

\newcommand*\diff{\mathop{}\!\mathrm{d}}
\newcommand*\sen{\operatorname{sen}}
\newcommand*\tg{\operatorname{tg}}
\newcommand*\arctg{\operatorname{arctg}}
\newcommand*\abs[1]{\left|#1\right|}

\newcommand*\tipo{PROVA IV}
\newcommand*\turma{MAT102-02U}
\newcommand*\disciplina{CDI0001}
\newcommand*\eu{Helder G. G. de Lima}
\newcommand*\data{07/12/2015}

\author{\eu}
\title{\tipo - \disciplina}
\date{\data}

\begin{document}
\thispagestyle{empty}
\newgeometry{margin=2cm,bottom=0.5cm}
\begin{center}
\includegraphics[width=9.0cm]{marca}
\noindent\begin{tabular}{l c c r}
  \textbf{\disciplina (\turma)}
& \textbf{\tipo}
& \textbf{\data}
\end{tabular}
\\ Prof. \eu\footnote{
Este é um material de acesso livre distribuído sob os termos da licença \href{https://creativecommons.org/licenses/by-sa/4.0/deed.pt_BR}{Creative Commons BY-SA 4.0}.}
\end{center}

\noindent Nome do(a) aluno(a): \rule{13cm}{0.01cm}

%\section*{Instruções}
\begin{center}\fbox{
\begin{minipage}{14cm}

{\footnotesize
\begin{itemize}
\renewcommand{\theenumi}{\Roman{enumi}}
\item Identifique-se em todas as folhas.
\item Mantenha o celular e os demais equipamentos eletrônicos desligados durante a prova.
\item Justifique cada resposta com cálculos ou argumentos baseados na teoria estudada.
\end{itemize}
}

\end{minipage}
}
\end{center}


%\section*{Questões}

\begin{enumerate}
\item Calcule apenas 5 das 6 integrais a seguir:
\begin{enumerate}
\item (2,0) $\displaystyle \int 5^x(\sen(5^x)-\cos(5^x)) \diff x$
\item (2,0) $\displaystyle \int 4x^3 \cos(x^2+1) \diff x$
\item (2,0) $\displaystyle \int \frac{8x+3}{x^2+3} \diff x$
\item (2,0) $\displaystyle \int e^{2x} \cos(x) \diff x$
\item (2,0) $\displaystyle \int \tg^3(x)\sec^3(x) \diff x$
\item (2,0) $\displaystyle \int \frac{2 x+13}{x^2-5 x+4} \diff x$
\end{enumerate}
\end{enumerate}

\begin{center}
BOA PROVA!!!
\end{center}
\begin{center}
(E BOAS FÉRIAS ;-)
\end{center}

\newpage
\section*{Respostas e observações}

\begin{enumerate}
\item Seguem alguns possíveis caminhos para o cálculo das integrais propostas:
\begin{enumerate}
\item $\displaystyle \int 5^x(\sen(5^x)-\cos(5^x)) \diff x$

Considerando $u = 5^x$, tem-se $\diff u = 5^x \ln(5) \diff x$ e
\begin{align*}
&   \int 5^x(\sen(5^x)-\cos(5^x)) \diff x
  = \frac{1}{\ln(5)} \int 5^x(\sen(5^x)-\cos(5^x)) \ln(5) \diff x\\
& = \frac{1}{\ln(5)} \int \sen(u) - \cos(u) \diff u
  = \frac{1}{\ln(5)} (-\cos(u) - \sen(u)) + k \\
& = \frac{-1}{\ln(5)} (\cos(5^x) + \sen(5^x)) + k.
\end{align*}

\item $\displaystyle \int 4x^3 \cos(x^2+1) \diff x$

Fazendo a substituição $t = x^2 + 1$, tem-se $x^2 = t-1$, $\diff t = 2x \diff x$ e
\begin{align*}
    \int 4x^3 \cos(x^2+1) \diff x
= 2 \int x^2 \cos(x^2+1) 2x \diff x
= 2 \int (t-1) \cos(t) \diff t
\end{align*}

Tomando $u = t-1$ e $\diff v = \cos(t) \diff t$, tem-se $\diff u = \diff t$ e $v = \sen(t)$. Logo:
\begin{align*}
  \int (t-1) \cos(t) \diff t
& = (t-1)\sen(t) - \int \sen(t) \diff t
  = (t-1)\sen(t) + \cos(t) + k_1.
\end{align*}

Portanto, $\int 4x^3 \cos(x^2+1) \diff x = 2 x^2 \sen(x^2+1) + 2\cos(x^2+1) + k_2$.

\item $\displaystyle \int \frac{8x+3}{x^2+3} \diff x$
Observe inicialmente que
\begin{align*}
  \int \frac{8x+3}{x^2+3} \diff x
= 4 \int \frac{2x}{x^2+3} \diff x + \int \frac{3}{x^2+3} \diff x.
\end{align*}
Para a primeira integral, é razoável escolher $u = x^2+3$, de modo que $\diff u = 2x \diff x$ e
\begin{align*}
  \int \frac{2x}{x^2+3} \diff x
= \int \frac{1}{u} \diff u
= \ln|u| + k_1
= \ln|x^2 + 3| + k_1
= \ln(x^2 + 3) + k_1
\end{align*}
Para a segunda, pode-se fazer a substituição trigonométrica $x = \sqrt{3} \tg(\theta)$ e obter $\diff x =\sqrt{3} \sec^2(\theta)\diff \theta$. Assim,
\begin{align*}
    \int \frac{3}{x^2+3} \diff x
& = \int \frac{3}{(\sqrt{3} \tg(\theta))^2+3} \sqrt{3} \sec^2(\theta)\diff \theta
  = \int \frac{3 \sqrt{3} \sec^2(\theta)}{3\tg^2(\theta)+3} \diff \theta \\
& = \sqrt{3} \int \frac{ \sec^2(\theta) }{ \tg^2(\theta) + 1 } \diff \theta
  = \sqrt{3} \int \frac{ \sec^2(\theta) }{ \sec^2(\theta)    } \diff \theta
  = \sqrt{3} \int 1 \diff \theta \\
& = \sqrt{3} \cdot \theta + k_2
  = \sqrt{3}\arctg\left(\frac{x}{\sqrt{3}}\right) + k_2
\end{align*}
Portanto, $\int \frac{8x+3}{x^2+3} \diff x
= \ln(x^2 + 3) + \sqrt{3}\arctg\left(\frac{x}{\sqrt{3}}\right) + k$.

\item $\displaystyle \int e^{2x} \cos(x) \diff x$

Pode-se aplicar a integração por partes com $u = \cos(x)$ e $\diff v = e^{2x} \diff x$, obtendo $\diff u = -\sen(x) \diff x$, $v = \frac{1}{2} e^{2x}$ e
\begin{align*}
  \int e^{2x} \cos(x) \diff x
= \frac{1}{2}\cos(x) e^{2x} + \frac{1}{2} \int e^{2x} \sen(x) \diff x.
\end{align*}
Novamente, é possível integrar por partes, considerando $u = \sen(x)$ e $\diff v = e^{2x} \diff x$, de modo que $\diff u = \cos(x) \diff x$, $v = \frac{1}{2} e^{2x}$ e
\begin{align*}
    \int e^{2x} \cos(x) \diff x
& = \frac{1}{2}\cos(x) e^{2x} + \frac{1}{2} \left[
\frac{1}{2}\sen(x) e^{2x} - \frac{1}{2} \int e^{2x} \cos(x) \diff x
\right] \\
& = \frac{1}{2}\cos(x) e^{2x}
  + \frac{1}{4}\sen(x) e^{2x}
  - \frac{1}{4}\int e^{2x} \cos(x) \diff x.
\end{align*}
Assim, somando $\frac{1}{4}\int e^{2x} \cos(x) \diff x$ a ambos os membros, resulta que
\begin{align*}
\frac{5}{4} \int e^{2x} \cos(x) \diff x
= \frac{1}{2}\cos(x) e^{2x}
+ \frac{1}{4}\sen(x) e^{2x} + k_1.
\end{align*}
Portanto,
\begin{align*}
  \int e^{2x} \cos(x) \diff x
= \frac{4}{5} \left[
    \frac{1}{2}\cos(x) e^{2x}
  + \frac{1}{4} \sen(x) e^{2x} + k_1
\right]
  = \frac{1}{5} e^{2x} (2 \cos(x) + \sen(x)) + k_2.
\end{align*}

\item $\displaystyle \int \tg^3(x)\sec^3(x) \diff x$

Usando a identidade trigonométrica $\tg^2(x) = \sec^2(x)-1$, segue que:
\begin{align*}
  \int \tg^3(x)\sec^3(x) \diff x
& = \int \tg^2(x)\tg(x)\sec^3(x) \diff x
  = \int \left[\sec^2(x) - 1\right]\tg(x)\sec^3(x) \diff x \\
& = \int \left[\sec^4(x) - \sec^2(x)\right]\tg(x)\sec(x) \diff x.
\end{align*}
Agora, basta substituir $u = \sec(x)$ para obter $\diff u = \tg(x)\sec(x) \diff x$ e
\begin{align*}
  \int \tg^3(x)\sec^3(x) \diff x
= \int u^4 - u^2 \diff u
= \frac{u^5}{5} - \frac{u^3}{3} + k
= \frac{\sec^5(x)}{5} - \frac{\sec^3(x)}{3} + k.
\end{align*}

\item $\displaystyle \int \frac{2x + 13}{x^2-5 x+4} \diff x$

Observando que $x^2 - 5x + 4 = (x-1)(x-4)$, a decomposição em frações parciais da função a ser integrada é obtida por
\[
\frac{2x + 13}{(x-1)(x-4)}
= \frac{A}{x-1} + \frac{B}{x-4}
= \frac{A(x-4) + B(x-1)}{(x-1)(x-4)}
= \frac{(A+B)x +(-4A-B)}{(x-1)(x-4)}.
\]

Comparando os numeradores, resulta que $A$ e $B$ devem ser soluções do sistema
\[
\begin{cases}
  A + B &= 2 \\
-4A - B &= 13
\end{cases}
\]

Somando termo a termo os dois membros das equações, resulta que $-3A = 15$, ou seja, $A=-5$. Consequentemente, $B = 2-(-5) = 7$. Então
\begin{align*}
\int \frac{2x + 13}{(x-1)(x-4)} \diff x
& = \int \frac{-5}{x-1} + \frac{7}{x-4} \diff x
  = -5 \int \frac{1}{x-1} \diff x + 7 \int \frac{1}{x-4} \diff x \\
& = -5 \ln\abs{ x-1 } + 7 \ln\abs{x-4} + c
  = \ln\abs{ (x-1)^{-5} (x-4)^7} + c.
\end{align*}

\end{enumerate}
\end{enumerate}

\end{document}
